\chapter{Ánh Sáng Từ Đông Sang Tây (Light from East to West)}
\begin{center}
  \textit{\color{truyengold}``Trí tuệ không có biên giới --- bài học của người xưa ở Đông hay Tây đều đến từ cùng một chỗ: trái tim con người.''}\\[4pt]
  \textit{(Wisdom knows no borders --- lessons of the ancients from East or West all come from the same place: the human heart.)}\\[4pt]
  \small\color{truyenblue}--- Cổ Kim Đông Tây ---
\end{center}
\ngancach

\section{Huyền Trang Thỉnh Kinh Tây Trúc}
\begin{truyen}{Huyền Trang Thỉnh Kinh Tây Trúc}{Đường Huyền Trang --- Trung Quốc, 629--645}
\chuhoa{N}{ăm} 629, nhà sư Huyền Trang một mình xuất phát từ Trường An đến Ấn Độ tìm kinh Phật, vượt qua sa mạc Gobi, núi Pamirs, và hàng ngàn dặm đường nguy hiểm.\\[4pt]
\textit{(In 629, the monk Xuanzang set out alone from Chang'an to India seeking Buddhist scriptures, crossing the Gobi Desert, the Pamir Mountains, and thousands of miles of dangerous road.)}

Ông không có bản đồ, không có đồng hành, không có sự bảo vệ của triều đình --- chỉ có niềm tin và ý chí.\\[4pt]
\textit{(He had no map, no companion, no imperial protection --- only faith and will.)}

Sau 16 năm và hơn 50.000 dặm, ông trở về với 657 bộ kinh sách --- và dành 19 năm còn lại của cuộc đời dịch những bộ kinh đó.\\[4pt]
\textit{(After 16 years and over 50,000 miles, he returned with 657 volumes of scriptures --- and devoted the remaining 19 years of his life to translating them.)}

Huyền Trang để lại cho đời không chỉ kinh sách --- mà còn bài học: kiến thức đáng để hy sinh cả cuộc đời để đi tìm.\\[4pt]
\textit{(Xuanzang left for the world not only scriptures --- but a lesson: knowledge is worth sacrificing a lifetime to seek.)}

\end{truyen}
\begin{baihoc}
  Hành trình tìm kiếm tri thức không bao giờ lãng phí --- dù con đường có dài bao nhiêu.\\[4pt]
  \textit{(The journey of seeking knowledge is never wasted --- no matter how long the road.)}
\end{baihoc}

\section{Ibn Battuta --- Nhà Lữ Hành Vĩ Đại Nhất Lịch Sử}
\begin{truyen}{Ibn Battuta --- Nhà Lữ Hành Vĩ Đại Nhất Lịch Sử}{Ibn Battuta --- Ma Rốc, 1325--1354}
\chuhoa{I}{bn} Battuta xuất phát từ Ma Rốc năm 21 tuổi để thực hiện cuộc hành hương Mecca --- và không dừng lại.\\[4pt]
\textit{(Ibn Battuta left Morocco at 21 to make the Hajj pilgrimage to Mecca --- and did not stop.)}

Trong 29 năm, ông đi qua 44 quốc gia ngày nay, tổng quãng đường 120.000 km --- vượt Marco Polo ba lần về khoảng cách.\\[4pt]
\textit{(Over 29 years, he traveled through what are today 44 countries, covering 120,000 km --- three times the distance of Marco Polo.)}

Ông học hỏi từ mỗi nền văn hóa, nói chuyện với học giả từ Trung Đông đến Trung Quốc, từ Tây Phi đến Ấn Độ.\\[4pt]
\textit{(He learned from every culture, conversed with scholars from the Middle East to China, from West Africa to India.)}

Sau này ông viết: ``Điều tôi học được từ mỗi nền văn hóa đều làm tôi trở nên người hơn, chứ không kém đi.''\\[4pt]
\textit{(Later he wrote: `What I learned from each culture made me more human, not less.')}

\end{truyen}
\begin{baihoc}
  Mỗi nền văn hóa đều có điều ta có thể học --- người khôn ngoan tìm tri thức ở khắp nơi, không chỉ trong những gì quen thuộc.\\[4pt]
  \textit{(Every culture has something we can learn --- the wise seek knowledge everywhere, not only in the familiar.)}
\end{baihoc}

\section{Marie Curie --- Phụ Nữ Đầu Tiên Đoạt Hai Giải Nobel}
\begin{truyen}{Marie Curie --- Phụ Nữ Đầu Tiên Đoạt Hai Giải Nobel}{Marie Curie --- Ba Lan/Pháp, 1867--1934}
\chuhoa{M}{arie} Curie sinh ra ở Ba Lan khi phụ nữ không được phép học đại học; bà phải bí mật tham gia một ``trường đại học di động'' bí mật của phong trào phụ nữ.\\[4pt]
\textit{(Marie Curie was born in Poland when women were not permitted to attend university; she secretly participated in a clandestine `floating university' of the women's movement.)}

Bà đến Paris học và phải chịu đựng cái lạnh cắt da, cái đói và sự kỳ thị --- và tốt nghiệp thủ khoa ngành Vật lý.\\[4pt]
\textit{(She went to Paris to study and endured bitter cold, hunger, and discrimination --- and graduated first in her class in Physics.)}

Bà đoạt giải Nobel Vật lý năm 1903 và Nobel Hóa học năm 1911 --- người duy nhất trong lịch sử đoạt hai giải Nobel ở hai ngành khoa học khác nhau.\\[4pt]
\textit{(She won the Nobel Prize in Physics in 1903 and Nobel Prize in Chemistry in 1911 --- the only person in history to win two Nobel Prizes in two different sciences.)}

Marie Curie nói: ``Trong cuộc sống không có gì đáng sợ, chỉ có những thứ cần được hiểu.''\\[4pt]
\textit{(Marie Curie said: `Nothing in life is to be feared, it is only to be understood.')}

\end{truyen}
\begin{baihoc}
  Không có rào cản nào cao hơn ý chí của người đủ kiên định để vượt qua --- kể cả rào cản xã hội và định kiến.\\[4pt]
  \textit{(No barrier is higher than the will of one determined enough to overcome it --- including social barriers and prejudice.)}
\end{baihoc}

\section{Gandhi Và Muối}
\begin{truyen}{Gandhi Và Muối}{Mahatma Gandhi --- Ấn Độ, 1930}
\chuhoa{N}{ăm} 1930, Gandhi dẫn 78 người đi bộ 388 km ra biển để tự làm muối --- một hành động vi phạm độc quyền muối của người Anh.\\[4pt]
\textit{(In 1930, Gandhi led 78 people on a 388-km walk to the sea to make their own salt --- an act defying the British salt monopoly.)}

Lúc đầu người Anh cười nhạo. Nhưng càng đi, đoàn người càng lớn --- hàng chục ngàn người nhập cuộc dọc đường.\\[4pt]
\textit{(At first the British laughed. But as the march continued, the procession grew --- tens of thousands joined along the way.)}

Hành động đơn giản là nhặt một nắm muối từ biển đã châm ngòi cho phong trào độc lập toàn Ấn Độ.\\[4pt]
\textit{(The simple act of lifting a handful of salt from the sea ignited the independence movement across all of India.)}

Gandhi dạy: ``Trước tiên họ bỏ qua bạn, rồi họ cười bạn, rồi họ chống lại bạn --- rồi bạn thắng.''\\[4pt]
\textit{(Gandhi taught: `First they ignore you, then they laugh at you, then they fight you --- then you win.')}

\end{truyen}
\begin{baihoc}
  Sức mạnh của chân lý và sự kiên trì không bạo lực --- được chứng minh bằng hành động đơn giản và kiên định --- có thể lay chuyển đế quốc.\\[4pt]
  \textit{(The power of truth and nonviolent persistence --- demonstrated through simple and steadfast action --- can shake empires.)}
\end{baihoc}
